% s6_mechanism.tex -- Section 6: pole structure, the closed-form approximation of the mode, and the shape of the law
% (main text; details in Supplementary Section S7).
% Numbers: every quantitative statement carries a comment  % src: <path relative to the root of the repository>.
% Tables are generated by code/article/s6_mechanism_tables.py -> data/article/s6_mechanism_tables.json
% (rows in data/article/s6_mechanism_rows.tex).  Figures: code/article/s3s6_modes_figures.py.
% Vocabulary: "closed-form approximation" is always Eq. (s6_mechanism:eq-L1); the maximiser of the two or three
% slowest exponentials with exact rates and residues is the "exact two-/three-exponential truncation".
\section{Pole structure, a closed-form approximation of the mode, and the shape of the law}\label{s6_mechanism:sec}

This section explains the results of Section~\ref{s5_modes:sec}; Section~\ref{s6b_rigorous:sec} states which parts of
the explanation are proved. We work with the continuous-time walk at unit rate, whose density is
$\gd(t)=\sum_j b_j\eu^{-\nu_j t}$ \eqref{s2_model:eq-spectral} and whose mean is $q\MF$; the discrete-time walk has the same
rates and residues. The propagator of the walk without target in the time domain is
$p(\vx,t\,|\,\vy)=\sum_{\vk}\phi_{\vk}(\vx)\phi_{\vk}(\vy)\eu^{-\mu_{\vk}t}$.

\subsection{Arrival and capture: a factorisation}\label{s6_mechanism:ssec-factor}

\begin{proposition}[Factorisation of the first-passage transform]\label{s6_mechanism:prop-factorisation}
Let $\vo\neq\va$, let $u(t)=\nn\,p(\va,t\,|\,\vo)$ be the occupation of the target site at time $t$ in the walk
without target, normalised by its stationary value $1/\nn$, and let $\FPT_{\mathrm u}$ be the first-passage time to
$\va$ when the start is drawn uniformly from $\Om$ (so that $\FPT_{\mathrm u}=0$ with probability $1/\nn$), with
transform $\Fhat_{\mathrm u}(s)=\E[\eu^{-s\FPT_{\mathrm u}}]$. Then, for $s>0$,
\begin{equation}\label{s6_mechanism:eq-factor}
  \Fhat_{\mathrm u}(s)=\frac{1}{\nn\,s\,\Phat(\va,s\,|\,\va)},\qquad
  \Fhat(s)=\bigl[s\,\hat u(s)\bigr]\,\Fhat_{\mathrm u}(s),
\end{equation}
that is, $\gd(t)=\int_{[0,t]}u'(t-\tau)\,\gd_{\mathrm u}(\dif\tau)$. For the corner start and the opposite corner target,
\begin{equation}\label{s6_mechanism:eq-u}
  u(t)=\Bigl[1+2\sum_{k=1}^{N-1}(-1)^{k}\,\ck{k}^{2}\,\eu^{-\varepsilon_k t}\Bigr]^{\dd},\qquad
  \varepsilon_k=\frac{2}{\dd}\,\sk{k}^{2},
\end{equation}
and for every fixed $x>0$, $u(x/\muone)\to\theta_4(\eu^{-x})^{\dd}$ as $N\to\infty$, where
$\theta_4(y)=1+2\sum_{k\ge1}(-1)^{k}y^{k^{2}}$ is the Jacobi theta function of the nome $y$ \citep[\S20.2]{DLMF}.
\end{proposition}

The proof (Supplementary Section~\ref{sup_mechanism:sec}) averages the renewal ratio over the start. The identity is the
renewal equation written differently, and we do not claim it as new; its use is that it separates the two time scales
of the problem. The factor $u$ describes how the walker \emph{arrives}: it knows nothing about the target and lives on
the diffusive scale $1/\muone\propto N^{2}$. The factor $\Fhat_{\mathrm u}$ describes \emph{capture} from equilibrium,
which for a reversible chain is close to exponential, with an error controlled by the ratio of the relaxation time to
the mean hitting time \citep{AldousBrown1992}, here of order $1/X$. In one dimension $X=2\sk{1}^{2}N(N-1)\to\pi^{2}/2$ is not
large and the two scales coincide; in two and three dimensions $X$ grows as $2\pi\ln N$ and as $7.107\,N$, so the target
is a weak sink on the diffusive scale.
% src: data/article/s6_mechanism_tables.json (asymptotes: X_const_2d = 0.38155, pi2_K3_over_6 = 7.10687)

\subsection{Formal two-pole asymptotics}\label{s6_mechanism:ssec-two-pole}

\begin{heuristic}[Closed-form approximation of the mode]\label{s6_mechanism:res-L1}
Let $\mu$ be the smallest non-zero relaxation rate (unit rate) of the reflecting lattice whose eigenspace has
non-zero weight at the target, $W=\nn\sum_{\mu_{\vk}=\mu}\phi_{\vk}(\va)^{2}$ that weight and
$A=-\nn\sum_{\mu_{\vk}=\mu}\phi_{\vk}(\va)\phi_{\vk}(\vo)$ the corresponding amplitude seen from the start, and let
$X=\mu\,q\MF$. If $A>0$ and $X\gg1$, the mode of the continuous-time first-passage density at unit rate is
\begin{equation}\label{s6_mechanism:eq-L1}
  \tmodec\;\approx\;q\MF\;\frac{\ln(A X)}{X+W-1},
\end{equation}
and the discrete-time mode at activity $q$ is $\tmode\approx\tmodec/q$. For corner-to-corner transport
$\mu=\muone=\tfrac{2}{\dd}\sin^{2}\tfrac{\pi}{2N}$ and $A=W=2\dd\cos^{2}\tfrac{\pi}{2N}$; replacing $\cos^{2}$ by 1 gives the
simplified form
\begin{equation}\label{s6_mechanism:eq-L1-simple}
  \frac{\tmode}{\MF}\;\approx\;\frac{\ln(2\dd X)}{X+2\dd-1},\qquad X=\muone\,q\MF .
\end{equation}
\end{heuristic}

\noindent We write $\tcf:=\MF\ln(2\dd X)/(X+2\dd-1)$ for the approximation of the discrete-time mode given by
\eqref{s6_mechanism:eq-L1-simple}; in Section~\ref{s6b_rigorous:sec} and in the accuracy statements of the Introduction
and of Section~\ref{s9_discussion:sec}, ``the closed form'' means $\tcf$, whereas ``the closed-form approximation''
\eqref{s6_mechanism:eq-L1} keeps the factors $A$ and $W$.

\noindent\emph{Derivation.} Separate the uniform eigenvector (rate $0$) and the eigenspace of $\mu$:
\begin{equation}\label{s6_mechanism:eq-split}
  \nn\,\Phat(\va,s\,|\,\va)=\frac1s+\frac{W}{s+\mu}+R(s),\qquad
  \hat u(s)=\nn\,\Phat(\va,s\,|\,\vo)=\frac1s-\frac{A}{s+\mu}+R_{\vo}(s),
\end{equation}
where $R$ and $R_{\vo}$ collect the eigenvalues $\mu_{\vk}>\mu$ and are analytic for $\operatorname{Re}s>-\mu'$, $\mu'$
being the next rate with weight at the target. The poles of $\Fhat$ are the zeros $-\nu_j$ of the first function, with
$0<\nu_0<\mu<\nu_1<\mu'$ (Proposition~\ref{s2_model:prop-interlacing}), and
$b_j=\hat u(-\nu_j)\big/\nn\,\partial_s\Phat(\va,s\,|\,\va)\big|_{s=-\nu_j}$. The limit $s\to0$ of $(1-\Fhat)/s$ gives
$q\MF=(W+A)/\mu+R(0)-R_{\vo}(0)$. \emph{Assume} that $X\gg1$, that $R$ varies by a relative amount $\Order(1/X)$ on
$[-\nu_1,0]$ with $R(0)=q\MF\,(1+\Order(1/X))$, and that $R_{\vo}=\Order(1/\mu)$ there. The two zero conditions then give
\begin{equation}\label{s6_mechanism:eq-poles}
  \nu_0=\frac{1}{q\MF}\bigl(1+\Order(X^{-1})\bigr),\qquad
  \nu_1=\mu\Bigl(1+\frac{W}{X}+\Order(X^{-2})\Bigr),\qquad
  b_0\approx\nu_0,\qquad b_1\approx-A\,\nu_0 .
\end{equation}
Keeping only $j=0,1$ in $\gd'(t)=-\sum_j b_j\nu_j\eu^{-\nu_j t}$ and solving $\gd'=0$ (possible because $b_1<0$, that is,
$A>0$) gives $\tmodec=\ln(-b_1\nu_1/b_0\nu_0)/(\nu_1-\nu_0)$. To leading order,
\begin{equation}\label{s6_mechanism:eq-L0}
  \tmodec\;\approx\;\frac{\ln(AX)}{\mu},
\end{equation}
and inserting the first-order rate difference $\nu_1-\nu_0=\mu\,(X+W-1)/X$ while keeping the argument of the logarithm at
its leading value $AX$ gives \eqref{s6_mechanism:eq-L1}.

\medskip\noindent\emph{Status.} This is a derivation, not a proof: it needs (i) uniform bounds on $R$ and $R'$, (ii) a bound
on the poles $j\ge2$, and (iii) unimodality. For corner-to-corner transport in two and three dimensions
Section~\ref{s6b_rigorous:sec} supplies all three by a route that does not expand $R$, and proves that
$X\muone(\tmodec-q\tcf)$ stays in a bounded window, a relative error $\Order(1/(X\ln X))$ for \eqref{s6_mechanism:eq-L1-simple}
(Theorem~\ref{s6b_rigorous:thm-cfcont}); the relations \eqref{s6_mechanism:eq-poles} are
Theorem~\ref{s6b_rigorous:thm-poles}(c). In one dimension the assumption $X\gg1$ fails and so does the formula (Refuted
statement~\ref{s6b_rigorous:ref-cf1d}). Equation~\eqref{s6_mechanism:eq-L1} is not a consistent expansion to order
$1/X$; the omitted terms partly cancel (Supplementary Section~\ref{sup_mechanism:ssec-ingredients}).

\subsection{Accuracy and leading laws}\label{s6_mechanism:ssec-accuracy}

Table~\ref{s6_mechanism:tab-accuracy} compares the predictions with the exact continuous-time modes, on the computed
grid of 41 sizes with $10\le N\le16\,777\,216$ for $\dd=2$ and 25 sizes with $10\le N\le4096$ for $\dd=3$.
% src: data/article/s6_mechanism_tables.json (accuracy/2d_CC, 3d_CC: n_sizes_N>=10); data/article/closed_form_checks.json
The leading-order form \eqref{s6_mechanism:eq-L0} errs by $+19.9\%$ at $N=10$ and still by $+2.4\%$ at $N=16\,777\,216$
for $\dd=2$. The closed-form approximation \eqref{s6_mechanism:eq-L1} is within $0.972\%$ of the exact mode for $\dd=2$
(worst at $N=113$) and within $0.289\%$ for $\dd=3$ (worst at $N=10$; within $1.3\times10^{-5}$ for $N\ge100$); the simplified
form \eqref{s6_mechanism:eq-L1-simple} is within $0.970\%$ and $0.085\%$. In two dimensions the error decreases slowly, to
$-0.45\%$ at $N=16\,777\,216$, where $X$ is only $104.9$. The errors of \eqref{s6_mechanism:eq-L1-simple} at the 40 sizes
with $N\ge12$ ($\dd=2$) and at all 25 sizes ($\dd=3$) lie inside the windows of Theorem~\ref{s6b_rigorous:thm-cfcont};
for $\dd=2$, $N=10$ no window is proved. With the two-term mean of Theorem~\ref{s4_mfpt:thm-expansion} and $X$ from
\eqref{s4_mfpt:eq-X2d}, so that no exact input is needed, the error of \eqref{s6_mechanism:eq-L1-simple} stays within
$0.975\%$ for $\dd=2$. The exact three-exponential truncation is within $5.35\times10^{-5}$ ($\dd=2$) and $7.1\times10^{-6}$
($\dd=3$): the mode is controlled by very few poles.
% src: data/article/closed_form_checks.json (2d_CC, 3d_CC: L1_weights, L1_plain, L0; 2d_CC_fully_asymptotic_L1_plain); data/article/s6_mechanism_tables.json (accuracy; accuracy/2d_CC/rows; three_pole); data/msc_modes/summary_tables.md (M2); data/article/s6b_rigorous_checks.json (continuous_time_against_windows: n_sizes_in_window_range, n_inside_window)

\begin{table}[tbp]
\centering
\caption{Signed relative error (in per cent, except the last column) of five predictions of the exact
continuous-time mode, corner to corner, unit rate: the leading-order form \eqref{s6_mechanism:eq-L0}; the
closed-form approximation \eqref{s6_mechanism:eq-L1} with $A=W=2\dd\cos^{2}(\pi/2N)$; its simplified form
\eqref{s6_mechanism:eq-L1-simple}; and the exact two- and three-exponential truncations, that is, the maximiser of
the two or three slowest exponentials with numerically exact rates and residues (last column: relative error, not
per cent). No parameter is fitted.}
\label{s6_mechanism:tab-accuracy}
% src: data/article/s6_mechanism_rows.tex, s6_mechanism_tables.json (accuracy/*/rows), computed by code/article/s6_mechanism_tables.py from data/msc_modes/laplace_modes.jsonl; L1 and L1s columns agree with data/article/closed_form_checks.json; L0, L1, two- and three-pole columns agree with data/msc_modes/summary_tables.md (M2, M3)
\small
\begin{tabular}{rrrrrrr}
\toprule
$N$ & $X$ & \eqref{s6_mechanism:eq-L0} (\%) & \eqref{s6_mechanism:eq-L1} (\%) & \eqref{s6_mechanism:eq-L1-simple} (\%) & two exp.\ (\%) & three exp.\\
\midrule
\multicolumn{7}{l}{$\dd=2$}\\
10 & 14.740 & $+19.939$ & $+0.2096$ & $+0.2660$ & $+0.6602$ & $+5.3\times10^{-5}$ \\
35 & 22.706 & $+12.231$ & $-0.8354$ & $-0.8222$ & $+0.5134$ & $+5.5\times10^{-6}$ \\
113 & 30.083 & $+8.901$ & $-0.9717$ & $-0.9700$ & $+0.3904$ & $+1.6\times10^{-6}$ \\
1280 & 45.335 & $+5.695$ & $-0.8654$ & $-0.8654$ & $+0.2437$ & $+3.1\times10^{-7}$ \\
10\,240 & 58.401 & $+4.361$ & $-0.7384$ & $-0.7384$ & $+0.1791$ & $+1.2\times10^{-7}$ \\
1\,048\,576 & 87.485 & $+2.876$ & $-0.5347$ & $-0.5347$ & $+0.1089$ & $+2.6\times10^{-8}$ \\
16\,777\,216 & 104.906 & $+2.392$ & $-0.4545$ & $-0.4545$ & $+0.0871$ & $+1.4\times10^{-8}$ \\
\midrule
\multicolumn{7}{l}{$\dd=3$}\\
10 & 64.168 & $+7.253$ & $-0.2886$ & $-0.0845$ & $+0.1544$ & $-7.1\times10^{-6}$ \\
35 & 242.186 & $+2.034$ & $-0.0253$ & $-0.0025$ & $+0.0365$ & $-3.1\times10^{-7}$ \\
100 & 704.235 & $+0.708$ & $-0.0013$ & $+0.0015$ & $+0.0112$ & $-2.9\times10^{-8}$ \\
320 & 2267.787 & $+0.221$ & $+0.0005$ & $+0.0008$ & $+0.0031$ & $-2.3\times10^{-9}$ \\
1280 & 9090.397 & $+0.055$ & $+0.0002$ & $+0.0002$ & $+0.0007$ & $-1.3\times10^{-10}$ \\
4096 & 29\,103.352 & $+0.017$ & $+0.0001$ & $+0.0001$ & $+0.0002$ & $+9.6\times10^{-11}$ \\
\bottomrule
\end{tabular}
\end{table}

With $1/\muone\simeq2\dd N^{2}/\pi^{2}$, $X\simeq2\pi\ln N$ for $\dd=2$ and $X\simeq\pi^{2}\Kthree N/6$ for $\dd=3$, the
leading-order form \eqref{s6_mechanism:eq-L0} becomes
\begin{equation}\label{s6_mechanism:eq-asym-2d}
  \dd=2:\qquad q\,\tmode\simeq\frac{4}{\pi^{2}}\,N^{2}\bigl[\ln\ln N+\ln 8\pi\bigr],\qquad
  \frac{\tmode}{\MF}\simeq\frac{\ln(8\pi\ln N)}{2\pi\ln N},
\end{equation}
\begin{equation}\label{s6_mechanism:eq-asym-3d}
  \dd=3:\qquad
  \begin{aligned}
  q\,\tmode&\simeq\frac{6}{\pi^{2}}\,N^{2}\bigl[\ln N+\ln(\pi^{2}\Kthree)\bigr]=0.6079\,N^{2}\,(\ln N+3.7528),\\
  \frac{\tmode}{\MF}&\simeq0.1407\,\frac{\ln N+3.753}{N}.
  \end{aligned}
\end{equation}
% src: data/article/s6_mechanism_tables.json (asymptotes: 6_over_pi2 = 0.60793, ln_pi2_K3 = 3.75282, 6/(pi^2 K3) = 0.14071)
Both laws are theorems with explicit remainders (Theorems~\ref{s6b_rigorous:thm-laws} and~\ref{s6b_rigorous:thm-disc}). In
three dimensions the law is visible in the data: it exceeds the exact $\tmodec/N^{2}$ by $8.7\%$ at $N=10$, $0.81\%$ at
$N=100$ and $0.02\%$ at $N=4096$. In two dimensions it is not: the constant $\ln8\pi=3.22$ exceeds $\ln\ln N$ for every
$N<\eu^{8\pi}\approx8\times10^{10}$, and the law exceeds the exact value by $19\%$, $8.9\%$ and $2.3\%$ at $N=10$, $100$ and
$16\,777\,216$. At every accessible size the two-dimensional mode should be described by
\eqref{s6_mechanism:eq-L1-simple}, not by its leading term.
% src: data/article/s6_mechanism_tables.json (asymptotes/3d: rel_err; asymptotes: ln_8pi, exp_8pi, 2d rel_err)

\subsection{Shape of the law: the mode is not a typical time}\label{s6_mechanism:ssec-shape}

Proposition~\ref{s6_mechanism:prop-factorisation} also describes the whole density. If $\gd_{\mathrm u}$ is replaced by the
exponential law of rate $\nu_0$, then $\gd(t)\approx\nu_0\int_0^{t}u'(\tau)\eu^{-\nu_0(t-\tau)}\dif\tau$, which is $\nu_0u(t)$ for
$t\ll q\MF$ and $\nu_0\eu^{-\nu_0t}$ for $t\gg1/\muone$, and whose maximum is where the late-time arrival rate
$u'(t)\approx A\mu\eu^{-\mu t}$ has dropped to the capture rate $\nu_0$, which is \eqref{s6_mechanism:eq-L0}.
Figure~\ref{s6_mechanism:fig-density} shows both regimes for the exact densities: in the variable $x=\muone t$ the product
$q\MF\,\gd$ follows $\theta_4(\eu^{-x})^{\dd}$ up to the maximum, and in the variable $t/(q\MF)$ it follows $\eu^{-t/(q\MF)}$.
There is no single-scale collapse.

\begin{figure}[tbp]
\centering
\includegraphics[width=\textwidth]{s6_mechanism_density.pdf}
\caption{Exact first-passage densities, corner to corner, continuous time at unit rate (lines; sizes and colours
as listed in panels (c) and (d)). (a,b) Against the diffusive time $x=\muone t$, for $\dd=2$ ($N=10$ to
$1\,048\,576$) and $\dd=3$ ($N=10$ to $4096$); dashed: $\theta_4(\eu^{-x})^{\dd}$; open circles: exact discrete-time
$\MF\,\pmf(t)$ at $q=0.8$ for $N=160$ ($\dd=2$) and $N=40$ ($\dd=3$), plotted against $\muone qt$. (c,d) The same
densities against $t/(q\MF)$ on a logarithmic scale; dashed: $\eu^{-t/(q\MF)}$.}
\label{s6_mechanism:fig-density}
% src: code/article/s3s6_modes_figures.py; data data/msc_modes/profiles.npz, data/msc_modes/profile_stats.json, data/msc_modes/pmf/
\end{figure}

\begin{table}[tbp]
\centering
\caption{Shape statistics of the exact first-passage law, corner to corner, continuous time at unit rate:
mode-to-mean ratio, probability of arrival by the modal time, median, survival at the mean, height of the
density at the mode in units of $1/q\MF$ (equal to 1 for an exponential law), and the window, in units of
$\tmodec$, in which the density exceeds $99\%$ of its maximum. For $\dd=1$ the start is the reflecting end.}
\label{s6_mechanism:tab-shape}
% src: data/article/s6_mechanism_rows.tex, s6_mechanism_tables.json (shape_table): columns 3-6 from data/msc_modes/medians.json; columns 7-8 recomputed by code/article/s6_mechanism_tables.py with code/msc_modes/fptlib.py (agree with data/msc_modes/profile_stats.json at the two sizes in common, differences below 4e-11)
\small
\begin{tabular}{rrcccccc}
\toprule
$\dd$ & $N$ & $\tmodec/q\MF$ & $\Prob(\FPT\le\tmodec)$ & median$/q\MF$ & $\Sv(q\MF)$ & $\gd(\tmodec)\,q\MF$ & $99\%$ window\\
\midrule
1 & 1000 & 0.3333 & 0.1665 & 0.7575 & 0.3708 & 0.9251 & $[0.89,\,1.13]$ \\
2 & 35 & 0.1769 & 0.0981 & 0.7195 & 0.3683 & 0.9448 & $[0.87,\,1.17]$ \\
2 & 200 & 0.1350 & 0.0780 & 0.7117 & 0.3681 & 0.9529 & $[0.85,\,1.19]$ \\
2 & 163\,840 & 0.0729 & 0.0456 & 0.7017 & 0.3679 & 0.9690 & $[0.82,\,1.27]$ \\
3 & 40 & 0.0262 & 0.0174 & 0.6959 & 0.3679 & 0.9879 & $[0.77,\,1.51]$ \\
3 & 100 & 0.0118 & 0.0082 & 0.6942 & 0.3679 & 0.9939 & $[0.72,\,1.97]$ \\
3 & 1280 & 0.0012 & 0.0009 & 0.6932 & 0.3679 & 0.9992 & $[0.58,\,9.47]$ \\
\bottomrule
\end{tabular}
\end{table}

Table~\ref{s6_mechanism:tab-shape} quantifies the consequence. For $\dd\ge2$ the first-passage law to a distant point
target is an exponential law preceded by a short dead time, not a peaked law with a long tail. Only $9.8\%$ ($\dd=2$,
$N=35$) and $1.7\%$ ($\dd=3$, $N=40$) of the walkers have arrived by the modal time, and this fraction decreases with $N$;
the median is $0.72\,\MF$ and $0.70\,\MF$ and approaches $\ln2=0.693$; the survival probability at the mean is $0.3683$ and
$0.3679$ against $1/\eu=0.3679$. The tail is exponential (geometric in discrete time) with rate $\nu_0$, not algebraic.
% src: data/msc_modes/medians.json; data/msc_modes/summary_tables.md (S, Md); data/article/s6_mechanism_tables.json (shape_table)
The maximum is also flat, increasingly so with dimension and size: for $\dd=3$ the density stays within $1\%$ of its
maximum from $0.58$ to $9.5$ modal times at $N=1280$, and at $\dd=3$, $N=100$ a relative error of $1\%$ in the density can
displace the apparent maximum anywhere between $0.72$ and $1.97$ modal times. The mode is therefore a soft statistic: it
marks the end of the dead time, it is not a typical arrival time, and the median is the appropriate typical value.
Statements about waiting times must be made with $\Sv$: the fraction of searches still unfinished is $36.8\%$ at
$t=\MF$, and $90\%$ ($\dd=2$, $N=35$) or $98\%$ ($\dd=3$, $N=40$) at $t=\tmode$.
% src: data/article/s6_mechanism_tables.json (shape_table, N = 1280); data/msc_modes/medians.json (S_at_mfpt 0.36827, 0.36788; S_at_mode 0.90189, 0.98261)

None of this structure is new in kind. Two time scales in confinement, and an exponential law rescaled by the mean
when the target is hard to find, are established for continuous and scale-invariant processes
\citep{Benichou2010,GodecMetzler2016}, and the rescaled laws fall into universality classes \citep{Meyer2011}; that the
mean may then be unrepresentative, and the most probable time far below it, is the subject of \citet{Mattos2012},
\citet{GodecMetzler2016} and \citet{Grebenkov2018}; and the exponential approximation for reversible chains is classical
\citep{AldousBrown1992}. In particular, \citet[Eq.~(31)]{Grebenkov2018} note that the median is close to $\ln2$ times the
mean when the survival probability is nearly exponential, which explains the limit $\ln2$ above, and give the
most probable time of the direct trajectories, $(r-\rho)^{2}/(6D)$ for a spherical target of radius $\rho$ at distance
$r$ \citep[Eq.~(3)]{Grebenkov2018}, the continuum counterpart of the first-image term of Section~\ref{s5_modes:sec}.
That estimate does not locate the mode of corner-to-corner transport for $\dd\ge2$: there the mode is set by the
relaxation of the box, $\tmode\approx\ln(AX)/(q\mu)$ for large $X$ by \eqref{s6_mechanism:eq-L1}, a time of order
$N^{2}\ln(AX)$ whose logarithmic factor grows with $N$, whereas the direct-trajectory time is of order $N^{2}$ with no
such factor. The relations
\eqref{s6_mechanism:eq-poles} are the lattice analogue of the perturbation of eigenvalues by a small removed region
\citep{WardKeller1993}, and uniform approximations of the density to a small target were given by
\citet{IsaacsonNewby2013}, who obtain the mode by maximising the density numerically, without a closed formula.
% src: Isaacson and Newby (2013), full text consulted in arXiv:1303.3016 (Eqs (2.11), (2.47)-(2.49); Sec. III.1; caption of Fig. 7: "we compute the mode by numerically maximizing the first passage time density")
% provenance of the attributions above: full texts read of IsaacsonNewby2013 (arXiv:1303.3016), GodecMetzler2016
% (arXiv:1611.07788), Benichou2010 (arXiv:1006.3477, rescaled FPT law exponential for non-compact exploration),
% Mattos2012 (arXiv:1206.1003) and Grebenkov2018 (arXiv:1911.00942, Eqs (3), (31)); Meyer2011 and WardKeller1993 are
% cited only for what their abstracts state (see the citedfor fields of refs.bib)
What is added here is the exact computation for the lattice walk, the dimension dependence of its mode, and the explicit
closed-form approximation \eqref{s6_mechanism:eq-L1} and closed form \eqref{s6_mechanism:eq-L1-simple} with their
measured accuracy; we do not know whether these formulas appear elsewhere and make no claim of priority.

\subsection{Other placements and the activity}\label{s6_mechanism:ssec-geometries}

For a corner start and a centre target (\CtM, odd $N$) the modes with $k_i=1$ vanish at the target, so that
$\mu=\tfrac{2}{\dd}\sin^{2}\tfrac{\pi}{N}\approx4\muone$, $W=2\dd$ and $A=2\dd\cos\tfrac{\pi}{N}$, and
\eqref{s6_mechanism:eq-L1} is within $0.964\%$ ($\dd=2$, 24 sizes, $11\le N\le655\,361$) and $0.206\%$ ($\dd=3$, 15
sizes, $11\le N\le1281$) of the exact mode. For a centre start and a corner target (\MtC) the formula does not apply:
the slowest mode seen by the target vanishes at the start ($A=0$), the residue $b_1$ is positive, and the maximum is set
by the first arrivals, to which several poles contribute; the mode is nevertheless far below the mean, with exact
discrete ratios $0.0352$ ($\dd=2$, $N=201$) and $0.0046$ ($\dd=3$, $N=61$) at $q=0.8$. For exit from the centre through an
absorbing boundary layer there is no small parameter, and $\tmode/\MF$ tends to a constant, $0.333284$, $0.474907$ and
$0.560839$ for $\dd=1,2,3$. Details are in Supplementary Section~\ref{sup_mechanism:ssec-geometries}.
% src: data/article/closed_form_checks.json (2d_C2M, 3d_C2M); data/article/s6_mechanism_tables.json (accuracy/2d_C2M, 3d_C2M; c2m_weights_check); data/msc_modes/discrete_modes.jsonl (M2C, q = 0.8: N = 201 ratio 0.035209; N = 61 ratio 0.004590); data/msc_modes/fit_results.json (continuum_exit)

By Proposition~\ref{s2_model:prop-structure}, $q\MF$ and $q$ times the continuous-time mode do not depend on $q$, and the
discrete-time mode inherits this up to a few steps: for $\dd=2$, $N=101$ the exact argmax times $q$ is $18\,056.1$,
$18\,056.0$, $18\,055.8$ and $18\,056$ for $q=0.3$, $0.5$, $0.9$ and $1$, against $\tmodec=18\,056.39$. Hence
$\tmode/\MF$ does not depend on $q$ beyond terms of order $1/\MF$.
% src: data/msc_modes/summary_tables.md (Q); data/msc_modes/discrete_modes.jsonl (group qscan); data/article/s6_mechanism_tables.json (activity)
